
\textbf{This defines the catala module \texttt{Faillissementspauliana}}

\section{Faillissementspauliana}



\subsection{Artikel 42 Faillissementswet}



\begin{tcolorbox}[colframe=OliveGreen, breakable, title=\textcolor{black}{\texttt{Metadata}},title after break=\textcolor{black}{\texttt{Metadata}},before skip=1em, after skip=1em, left=0.6em, right=0.6em]
\begin{Verbatim}[commandchars=\\\{\},numbers=left,firstnumber=8,stepnumber=1,label={\hspace*{\fill}\texttt{faillissementspauliana.catala\_en}},numbersep=9mm]
```catala
\PY{k+kr}{declaration} \PY{k+kt}{structure} \PY{n+nc}{Handeling}\PY{o}{:}
\PY{c+c1}{  \PYZsh{} Handeling kan ook feitelijke handeling zijn, valt dan buiten scope}
\PY{c+c1}{  \PYZsh{} Zie in dit kader ook HR 21 juni 2013 (Eringa q.q./ABN AMRO)}
  \PY{k+kr}{data} \PY{n+nv}{is\PYZus{}rechtshandeling} \PY{k+kr}{content} \PY{k+kt}{boolean}
  \PY{k+kr}{data} \PY{n+nv}{datum\PYZus{}rechtshandeling} \PY{k+kr}{content} \PY{k+kt}{date}
  \PY{k+kr}{data} \PY{n+nv}{verricht\PYZus{}door\PYZus{}de\PYZus{}schuldenaar} \PY{k+kr}{content} \PY{k+kt}{boolean}
  \PY{k+kr}{data} \PY{n+nv}{onverplicht} \PY{k+kr}{content} \PY{k+kt}{boolean}
\PY{c+c1}{  \PYZsh{} \PYZdq{}anders dan om niet\PYZdq{} (om baat) wordt gemodelleerd als negatie: not om\PYZus{}niet}
  \PY{k+kr}{data} \PY{n+nv}{om\PYZus{}niet} \PY{k+kr}{content} \PY{k+kt}{boolean}

\PY{k+kr}{declaration} \PY{k+kt}{structure} \PY{n+nc}{Faillissement}\PY{o}{:}
  \PY{k+kr}{data} \PY{n+nv}{datum\PYZus{}faillietverklaring} \PY{k+kr}{content} \PY{k+kt}{date}

\PY{k+kr}{declaration} \PY{k+kr}{scope} \PY{n+nc}{FaillissementspaulianaArt42}\PY{o}{:}
  \PY{k+kd}{input}\PY{+w}{ }\PY{n+nv}{handeling} \PY{k+kr}{content} \PY{n+nc}{Handeling}
  \PY{k+kd}{input}\PY{+w}{ }\PY{n+nv}{faillissement} \PY{k+kr}{content} \PY{n+nc}{Faillissement}

  \PY{k+kd}{input}\PY{+w}{ }\PY{n+nv}{benadeling\PYZus{}schuldeisers} \PY{k+kr}{condition}
  \PY{k+kd}{input}\PY{+w}{ }\PY{n+nv}{wetenschap\PYZus{}schuldenaar\PYZus{}aangetoond} \PY{k+kr}{condition}
\PY{c+c1}{  \PYZsh{} Schuldenaar\PYZhy{}wetenschap komt bij elke redenering terug maar bij de}
\PY{c+c1}{  \PYZsh{} wederpartij niet altijd.}
  \PY{k+kd}{context}\PY{+w}{ }\PY{n+nv}{wetenschap\PYZus{}wederpartij\PYZus{}aangetoond} \PY{k+kr}{condition}

\PY{c+c1}{  \PYZsh{} Hieronder wordt context gebruikt omdat deze dan niet voor iedere redenering}
\PY{c+c1}{  \PYZsh{} verplicht zijn om in te vullen.}

\PY{c+c1}{  \PYZsh{} \PYZdq{}behoudens tegenbewijs\PYZdq{}: geslaagd tegenbewijs ontkracht het bewijsvermoeden,}
\PY{c+c1}{  \PYZsh{} dan valt de wetenschap terug op de hoofdregel van art. 42 (de curator draagt}
\PY{c+c1}{  \PYZsh{} het bewijsrisico). Per zijde 1 conditie: art. 43 vermoedt aan beide}
\PY{c+c1}{  \PYZsh{} zijden, art. 45 enkel bij de schuldenaar.}
  \PY{k+kd}{context}\PY{+w}{ }\PY{n+nv}{vermoeden\PYZus{}wetenschap\PYZus{}schuldenaar\PYZus{}tegenbewijs} \PY{k+kr}{condition}
  \PY{k+kd}{context}\PY{+w}{ }\PY{n+nv}{vermoeden\PYZus{}wetenschap\PYZus{}wederpartij\PYZus{}tegenbewijs} \PY{k+kr}{condition}

\PY{c+c1}{  \PYZsh{} Gronden (\PYZdq{}verdachte categorieën\PYZdq{}) van het bewijsvermoeden van art. 43 lid 1.}
  \PY{k+kd}{context}\PY{+w}{ }\PY{n+nv}{tegen\PYZus{}onderwaarde} \PY{k+kr}{condition}
  \PY{k+kd}{context}\PY{+w}{ }\PY{n+nv}{voldoening\PYZus{}of\PYZus{}zekerheid\PYZus{}niet\PYZus{}opeisbare\PYZus{}schuld} \PY{k+kr}{condition}
  \PY{k+kd}{context}\PY{+w}{ }\PY{n+nv}{gelieerde\PYZus{}partij} \PY{k+kr}{condition}
\PY{c+c1}{  \PYZsh{} Art. 43 lid 1 aanvullende vereiste: de schuldenaar had zich niet reeds voor}
\PY{c+c1}{  \PYZsh{} de aanvang van de jaartermijn tot de handeling verplicht.}
  \PY{k+kd}{context}\PY{+w}{ }\PY{n+nv}{reeds\PYZus{}verplicht\PYZus{}voor\PYZus{}aanvang\PYZus{}termijn} \PY{k+kr}{condition}

\PY{c+c1}{  \PYZsh{} Bewijsvermoeden bij een rechtshandeling om niet (art. 45).}
  \PY{n+nv}{bewijsvermoeden\PYZus{}art45} \PY{k+kr}{scope} \PY{n+nc}{BewijsvermoedenArt45}

\PY{c+c1}{  \PYZsh{} Bewijsvermoeden van art. 43 (voornamelijk om baat).}
  \PY{n+nv}{bewijsvermoeden\PYZus{}art43} \PY{k+kr}{scope} \PY{n+nc}{BewijsvermoedenArt43}

  \PY{k+kd}{internal}\PY{+w}{ }\PY{n+nv}{voor\PYZus{}faillietverklaring} \PY{k+kr}{condition}
  \PY{k+kd}{internal}\PY{+w}{ }\PY{n+nv}{wetenschap\PYZus{}schuldenaar} \PY{k+kr}{condition}
  \PY{k+kd}{internal}\PY{+w}{ }\PY{n+nv}{wetenschap\PYZus{}wederpartij} \PY{k+kr}{condition}
  \PY{k+kd}{internal}\PY{+w}{ }\PY{n+nv}{wetenschapsvereiste} \PY{k+kr}{condition}

  \PY{k+kd}{output}\PY{+w}{ }\PY{n+nv}{vernietigbaar} \PY{k+kr}{condition}
```
\end{Verbatim}

\end{tcolorbox}



\subsubsection{Lid 1}



De curator kan ten behoeve van de boedel elke rechtshandeling die de
schuldenaar vóór de faillietverklaring onverplicht heeft verricht en
waarvan deze bij dit verrichten wist of behoorde te weten dat daarvan
benadeling van de schuldeisers het gevolg zou zijn, door een
buitengerechtelijke verklaring vernietigen. Artikel 50, tweede lid, van
Boek 3 van het Burgerlijk Wetboek is niet van toepassing.



\begin{Verbatim}[commandchars=\\\{\},numbers=left,firstnumber=68,stepnumber=1,label={\hspace*{\fill}\texttt{faillissementspauliana.catala\_en}}]
```catala
\PY{k+kr}{scope} \PY{n+nc}{FaillissementspaulianaArt42}\PY{o}{:}
\PY{c+c1}{  \PYZsh{} Fixatiebeginsel (art. 23 Fw)}
  \PY{k+kr}{rule} \PY{n+nv}{voor\PYZus{}faillietverklaring}
  \PY{k+kr}{under condition}
    \PY{n+nv}{handeling}\PY{o}{.}datum\PYZus{}rechtshandeling \PY{o}{\PYZlt{}} \PY{n+nv}{faillissement}\PY{o}{.}datum\PYZus{}faillietverklaring
  \PY{k+kr}{consequence} \PY{k+kr}{fulfilled}

  \PY{k+kr}{rule} \PY{n+nv}{wetenschap\PYZus{}schuldenaar}
  \PY{k+kr}{under condition} \PY{n+nv}{wetenschap\PYZus{}schuldenaar\PYZus{}aangetoond}
  \PY{k+kr}{consequence} \PY{k+kr}{fulfilled}

\PY{c+c1}{  \PYZsh{} Lid 1 verlangt steeds wetenschap bij de schuldenaar. Lid 2 scherpt dit}
\PY{c+c1}{  \PYZsh{} aan voor rechtshandelingen anders dan om niet.}
  \PY{k+kr}{rule} \PY{n+nv}{wetenschapsvereiste}
  \PY{k+kr}{under condition} \PY{n+nv}{wetenschap\PYZus{}schuldenaar}
  \PY{k+kr}{consequence} \PY{k+kr}{fulfilled}

\PY{c+c1}{  \PYZsh{} De rechtshandeling is vernietigbaar indien aan alle vereisten van lid 1}
\PY{c+c1}{  \PYZsh{} is voldaan.}
  \PY{k+kr}{rule} \PY{n+nv}{vernietigbaar}
  \PY{k+kr}{under condition}
    \PY{n+nv}{handeling}\PY{o}{.}is\PYZus{}rechtshandeling
    \PY{o}{and} \PY{n+nv}{handeling}\PY{o}{.}verricht\PYZus{}door\PYZus{}de\PYZus{}schuldenaar
    \PY{o}{and} \PY{n+nv}{voor\PYZus{}faillietverklaring}
    \PY{o}{and} \PY{n+nv}{handeling}\PY{o}{.}onverplicht
    \PY{o}{and} \PY{n+nv}{benadeling\PYZus{}schuldeisers}
    \PY{o}{and} \PY{n+nv}{wetenschapsvereiste}
  \PY{k+kr}{consequence} \PY{k+kr}{fulfilled}
```
\end{Verbatim}




\subsubsection{Lid 2}



Een rechtshandeling anders dan om niet, die hetzij meerzijdig is, hetzij
eenzijdig en tot een of meer bepaalde personen gericht, kan wegens
benadeling slechts worden vernietigd, indien ook degenen met of jegens
wie de schuldenaar de rechtshandeling verrichtte, wisten of behoorden te
weten dat daarvan benadeling van de schuldeisers het gevolg zou zijn.



\begin{Verbatim}[commandchars=\\\{\},numbers=left,firstnumber=103,stepnumber=1,label={\hspace*{\fill}\texttt{faillissementspauliana.catala\_en}}]
```catala
\PY{k+kr}{scope} \PY{n+nc}{FaillissementspaulianaArt42}\PY{o}{:}
\PY{c+c1}{  \PYZsh{} Wetenschapsvereiste wederpartij vereist bij \PYZdq{}anders dan om niet\PYZdq{}.}
  \PY{k+kr}{rule} \PY{n+nv}{wetenschap\PYZus{}wederpartij}
  \PY{k+kr}{under condition} \PY{n+nv}{wetenschap\PYZus{}wederpartij\PYZus{}aangetoond}
  \PY{k+kr}{consequence} \PY{k+kr}{fulfilled}

\PY{c+c1}{  \PYZsh{} Bij een rechtshandeling om niet is wetenschap bij wederpartij niet vereist.}
  \PY{k+kr}{exception}
  \PY{k+kr}{rule} \PY{n+nv}{wetenschapsvereiste}
  \PY{k+kr}{under condition}
    \PY{o}{(}\PY{o}{not} \PY{n+nv}{handeling}\PY{o}{.}om\PYZus{}niet\PY{o}{)}
    \PY{o}{and} \PY{o}{(}\PY{o}{not} \PY{n+nv}{wetenschap\PYZus{}wederpartij}\PY{o}{)}
  \PY{k+kr}{consequence} \PY{o}{not} \PY{k+kr}{fulfilled}
```
\end{Verbatim}




\subsubsection{Lid 3}



Wordt een rechtshandeling om niet wegens benadeling vernietigd, dan
heeft de vernietiging ten aanzien van de bevoordeelde, die wist noch
behoorde te weten dat van de rechtshandeling benadeling van de
schuldeisers het gevolg zou zijn, geen werking, voorzover hij aantoont
dat hij ten tijde van de faillietverklaring niet ten gevolge van de
rechtshandeling gebaat was.



\begin{Verbatim}[commandchars=\\\{\},numbers=left,firstnumber=123,stepnumber=1,label={\hspace*{\fill}\texttt{faillissementspauliana.catala\_en}}]
```catala
\PY{c+c1}{\PYZsh{} Bij rechtshandelingen om niet geldt de wetenschapseis aan de zijde van de}
\PY{c+c1}{\PYZsh{} wederpartij niet. Dit is al bij lid 1 gemodelleerd.}
\PY{c+c1}{\PYZsh{} Bescherming van derden (lid 3) is voor nu bewust buiten het model gelaten.}
```
\end{Verbatim}




\subsection{Artikel 43 Faillissementswet}



\begin{tcolorbox}[colframe=OliveGreen, breakable, title=\textcolor{black}{\texttt{Metadata}},title after break=\textcolor{black}{\texttt{Metadata}},before skip=1em, after skip=1em, left=0.6em, right=0.6em]
\begin{Verbatim}[commandchars=\\\{\},numbers=left,firstnumber=131,stepnumber=1,label={\hspace*{\fill}\texttt{faillissementspauliana.catala\_en}},numbersep=9mm]
```catala
\PY{k+kr}{declaration} \PY{k+kr}{scope} \PY{n+nc}{BewijsvermoedenArt43}\PY{o}{:}
  \PY{k+kd}{input}\PY{+w}{ }\PY{n+nv}{handeling} \PY{k+kr}{content} \PY{n+nc}{Handeling}
  \PY{k+kd}{input}\PY{+w}{ }\PY{n+nv}{faillissement} \PY{k+kr}{content} \PY{n+nc}{Faillissement}

  \PY{k+kd}{input}\PY{+w}{ }\PY{n+nv}{tegen\PYZus{}onderwaarde} \PY{k+kr}{condition}
  \PY{k+kd}{input}\PY{+w}{ }\PY{n+nv}{voldoening\PYZus{}of\PYZus{}zekerheid\PYZus{}niet\PYZus{}opeisbare\PYZus{}schuld} \PY{k+kr}{condition}
\PY{c+c1}{  \PYZsh{} Derde verdachte categorie gelieerde\PYZus{}partij. Sub 3\PYZhy{}6 wijzen, met de}
\PY{c+c1}{  \PYZsh{} begripsbepalingen van lid 2\PYZhy{}5, een reeks gelieerde partijen aan. Deze}
\PY{c+c1}{  \PYZsh{} gronden zijn sterk talig, ze worden samengevat in 1 conditie in plaats van}
\PY{c+c1}{  \PYZsh{} in een groot aantal afzonderlijke booleans uitgewerkt. Het verder modelleren}
\PY{c+c1}{  \PYZsh{} zou de neurale verwerking mogelijk met veel context belasten}
\PY{c+c1}{  \PYZsh{} (\PYZdq{}information overload\PYZdq{}, zie ook RAG\PYZhy{}survey Gao et al. 2024).}
\PY{c+c1}{  \PYZsh{} De kwalificatie als gelieerde partij wordt daarom neuraal bepaald.}
  \PY{k+kd}{input}\PY{+w}{ }\PY{n+nv}{gelieerde\PYZus{}partij} \PY{k+kr}{condition}
  \PY{k+kd}{input}\PY{+w}{ }\PY{n+nv}{reeds\PYZus{}verplicht\PYZus{}voor\PYZus{}aanvang\PYZus{}termijn} \PY{k+kr}{condition}

\PY{c+c1}{  \PYZsh{} Per zijde één weerlegbaar vermoeden (\PYZdq{}vermoed aan beide zijden te bestaan,}
\PY{c+c1}{  \PYZsh{} behoudens tegenbewijs\PYZdq{}). Ontzenuwen van bewijsvermoeden aan beide zijden}
\PY{c+c1}{  \PYZsh{} lossen we neuraal op (beide condities op False).}
  \PY{k+kd}{input}\PY{+w}{ }\PY{n+nv}{vermoeden\PYZus{}wetenschap\PYZus{}schuldenaar\PYZus{}tegenbewijs} \PY{k+kr}{condition}
  \PY{k+kd}{input}\PY{+w}{ }\PY{n+nv}{vermoeden\PYZus{}wetenschap\PYZus{}wederpartij\PYZus{}tegenbewijs} \PY{k+kr}{condition}

  \PY{k+kd}{internal}\PY{+w}{ }\PY{n+nv}{binnen\PYZus{}een\PYZus{}jaar\PYZus{}voor\PYZus{}faillietverklaring} \PY{k+kr}{condition}
  \PY{k+kd}{internal}\PY{+w}{ }\PY{n+nv}{niet\PYZus{}reeds\PYZus{}verplicht\PYZus{}voor\PYZus{}aanvang\PYZus{}termijn} \PY{k+kr}{condition}
  \PY{k+kd}{internal}\PY{+w}{ }\PY{n+nv}{verdachte\PYZus{}categorie} \PY{k+kr}{condition}
  \PY{k+kd}{internal}\PY{+w}{ }\PY{n+nv}{vermoeden\PYZus{}van\PYZus{}toepassing} \PY{k+kr}{condition}

  \PY{k+kd}{output}\PY{+w}{ }\PY{n+nv}{vermoeden\PYZus{}wetenschap\PYZus{}schuldenaar} \PY{k+kr}{condition}
  \PY{k+kd}{output}\PY{+w}{ }\PY{n+nv}{vermoeden\PYZus{}wetenschap\PYZus{}wederpartij} \PY{k+kr}{condition}
```
\end{Verbatim}

\end{tcolorbox}



\subsubsection{Lid 1}



Indien de rechtshandeling waardoor de schuldeisers zijn benadeeld, is
verricht binnen een jaar voor de faillietverklaring en de schuldenaar
zich niet reeds voor de aanvang van die termijn daartoe had verplicht,
wordt de aan het slot van artikel 42, eerste lid, eerste zin, bedoelde
wetenschap, behoudens tegenbewijs, vermoed aan beide zijden te bestaan:



\subsubsubsection{Sub 1°}



bij overeenkomsten, waarbij de waarde der verbintenis aan de zijde van
de schuldenaar aanmerkelijk die der verbintenis aan de andere zijde
overtreft;



\subsubsubsection{Sub 2°}



bij rechtshandelingen ter voldoening van of zekerheidstelling voor een
niet opeisbare schuld;



\subsubsubsection{Sub 3°}



bij rechtshandelingen, door de schuldenaar die een natuurlijk persoon
is, verricht met of jegens: a) zijn echtgenoot, zijn pleegkind of een
bloed- of aanverwant tot in de derde graad; b) een rechtspersoon waarin
hij, zijn echtgenoot, zijn pleegkind of een bloed- of aanverwant tot in
de derde graad bestuurder of commissaris is, dan wel waarin deze
personen, afzonderlijk of tezamen, als aandeelhouder rechtstreeks of
middellijk voor ten minste de helft van het geplaatste kapitaal
deelnemen;



\subsubsubsection{Sub 4°}



bij rechtshandelingen, door de schuldenaar die rechtspersoon is,
verricht met of jegens een natuurlijk persoon, a) die bestuurder of
commissaris van de rechtspersoon is, dan wel met of jegens diens
echtgenoot, pleegkind of bloed- of aanverwant tot in de derde graad: b)
die al dan niet tezamen met zijn echtgenoot, zijn pleegkinderen en zijn
bloed- of aanverwanten tot in de derde graad, als aandeelhouder
rechtstreeks of middellijk voor ten minste de helft van het geplaatste
kapitaal deelneemt; c) wiens echtgenoot, pleegkinderen of bloed- of
aanverwanten tot in de derde graad, afzonderlijk of tezamen, als
aandeelhouder rechtstreeks of middellijk voor tenminste de helft van het
geplaatste kapitaal deelnemen;



\subsubsubsection{Sub 5°}



bij rechtshandelingen, door de schuldenaar die rechtspersoon is,
verricht met of jegens een andere rechtspersoon, indien a) een van deze
rechtspersonen bestuurder is van de andere; b) een bestuurder,
natuurlijk persoon, van een van deze rechtspersonen, of diens
echtgenoot, pleegkind of bloed- of aanverwant tot in de derde graad,
bestuurder is van de andere; c) een bestuurder, natuurlijk persoon, of
een commissaris van een van deze rechtspersonen, of diens echtgenoot,
pleegkind of bloed- of aanverwant tot in de derde graad, afzonderlijk of
tezamen, als aandeelhouder rechtstreeks of middellijk voor ten minste de
helft van het geplaatste kapitaal deelneemt in de andere; d) in beide
rechtspersonen voor ten minste de helft van het geplaatste kapitaal
rechtstreeks of middellijk wordt deelgenomen door dezelfde
rechtspersoon, dan wel dezelfde natuurlijke persoon, al dan niet tezamen
met zijn echtgenoot, zijn pleegkinderen en zijn bloed- of aanverwanten
tot in de derde graad;



\subsubsubsection{Sub 6°}



bij rechtshandelingen, door de schuldenaar die rechtspersoon is,
verricht met of jegens een groepsmaatschappij.



\subsubsection{Lid 2}



Met een echtgenoot wordt een geregistreerde partner of een andere
levensgezel gelijkgesteld.



\subsubsubsection{Lid 3}



Onder pleegkind wordt verstaan hij die duurzaam als eigen kind is
verzorgd en opgevoed.



\subsubsubsection{Lid 4}



Onder bestuurder, commissaris of aandeelhouder wordt mede verstaan hij
die minder dan een jaar vóór de rechtshandeling bestuurder, commissaris
of aandeelhouder is geweest.



\subsubsubsection{Lid 5}



Indien de bestuurder van een rechtspersoon-bestuurder zelf een
rechtspersoon is, wordt deze rechtspersoon met de
rechtspersoon-bestuurder gelijkgesteld.



\subsubsubsection{Lid 6}



Artikel 138, tiende lid, van boek 2 van het Burgerlijk Wetboek is van
toepassing ingeval de schuldenaar een rechtspersoon is.



\begin{Verbatim}[commandchars=\\\{\},numbers=left,firstnumber=220,stepnumber=1,label={\hspace*{\fill}\texttt{faillissementspauliana.catala\_en}}]
```catala
\PY{k+kr}{scope} \PY{n+nc}{BewijsvermoedenArt43}\PY{o}{:}
\PY{c+c1}{  \PYZsh{} Rekenen met data in een schrikkeljaar is complex (29 febr. in een}
\PY{c+c1}{  \PYZsh{} schrikkeljaar minus 1 jaar?)}
\PY{c+c1}{  \PYZsh{} De wetgever heeft zich hier niet duidelijk over uitgelaten tav de Fw.}
\PY{c+c1}{  \PYZsh{} Voor nu is de optiek van de HR in het burgerlijk procesrecht overgenomen}
\PY{c+c1}{  \PYZsh{} voor het datumrekenen. Zie HR 1 september 2017, ECLI:NL:HR:2017:2225}
  \PY{k+kt}{date} \PY{k+kd}{round} \PY{k+kr}{down}

\PY{c+c1}{  \PYZsh{} \PYZdq{}verricht binnen een jaar voor de faillietverklaring\PYZdq{}}
  \PY{k+kr}{rule} \PY{n+nv}{binnen\PYZus{}een\PYZus{}jaar\PYZus{}voor\PYZus{}faillietverklaring}
  \PY{k+kr}{under condition}
    \PY{o}{(}\PY{n+nv}{handeling}\PY{o}{.}datum\PYZus{}rechtshandeling
      \PY{o}{\PYZgt{}}\PY{o}{=} \PY{n+nv}{faillissement}\PY{o}{.}datum\PYZus{}faillietverklaring \PY{o}{\PYZhy{}} \PY{l+m+mi}{1} \PY{o}{year}\PY{o}{)}
    \PY{o}{and} \PY{o}{(}\PY{n+nv}{handeling}\PY{o}{.}datum\PYZus{}rechtshandeling
      \PY{o}{\PYZlt{}} \PY{n+nv}{faillissement}\PY{o}{.}datum\PYZus{}faillietverklaring\PY{o}{)}
  \PY{k+kr}{consequence} \PY{k+kr}{fulfilled}

\PY{c+c1}{  \PYZsh{} \PYZdq{}en de schuldenaar zich niet reeds voor de aanvang van die termijn daartoe}
\PY{c+c1}{  \PYZsh{} had verplicht\PYZdq{}}
  \PY{k+kr}{rule} \PY{n+nv}{niet\PYZus{}reeds\PYZus{}verplicht\PYZus{}voor\PYZus{}aanvang\PYZus{}termijn}
  \PY{k+kr}{under condition} \PY{o}{not} \PY{n+nv}{reeds\PYZus{}verplicht\PYZus{}voor\PYZus{}aanvang\PYZus{}termijn}
  \PY{k+kr}{consequence} \PY{k+kr}{fulfilled}

\PY{c+c1}{  \PYZsh{} OF\PYZhy{}OF: het volstaat dat de handeling onder één van de drie categorieën valt.}
  \PY{k+kr}{rule} \PY{n+nv}{verdachte\PYZus{}categorie}
  \PY{k+kr}{under condition}
    \PY{n+nv}{tegen\PYZus{}onderwaarde}
    \PY{o}{or} \PY{n+nv}{voldoening\PYZus{}of\PYZus{}zekerheid\PYZus{}niet\PYZus{}opeisbare\PYZus{}schuld}
    \PY{o}{or} \PY{n+nv}{gelieerde\PYZus{}partij}
  \PY{k+kr}{consequence} \PY{k+kr}{fulfilled}

\PY{c+c1}{  \PYZsh{} Anders dan bij art. 45 is hier geen `not om\PYZus{}niet` vereist: ook bij een}
\PY{c+c1}{  \PYZsh{} rechtshandeling om niet kunnen de vermoedens van art. 43 worden ingeroepen}
\PY{c+c1}{  \PYZsh{} zie in dit kader r.o. 3.5.7 \PYZhy{} 3.5.8 van HR 23 september 2016,}
\PY{c+c1}{  \PYZsh{} ECLI:NL:HR:2016:2172 (Hage q.q./Kelderman).}
  \PY{k+kr}{rule} \PY{n+nv}{vermoeden\PYZus{}van\PYZus{}toepassing}
  \PY{k+kr}{under condition}
    \PY{n+nv}{binnen\PYZus{}een\PYZus{}jaar\PYZus{}voor\PYZus{}faillietverklaring}
    \PY{o}{and} \PY{n+nv}{niet\PYZus{}reeds\PYZus{}verplicht\PYZus{}voor\PYZus{}aanvang\PYZus{}termijn}
    \PY{o}{and} \PY{n+nv}{verdachte\PYZus{}categorie}
  \PY{k+kr}{consequence} \PY{k+kr}{fulfilled}

\PY{c+c1}{  \PYZsh{} geslaagd tegenbewijs onthoudt enkel die zijde het vermoeden, waarna de}
\PY{c+c1}{  \PYZsh{} wetenschap terugvalt op de hoofdregel van art. 42.}
  \PY{k+kr}{rule} \PY{n+nv}{vermoeden\PYZus{}wetenschap\PYZus{}schuldenaar}
  \PY{k+kr}{under condition}
    \PY{n+nv}{vermoeden\PYZus{}van\PYZus{}toepassing}
    \PY{o}{and} \PY{o}{not} \PY{n+nv}{vermoeden\PYZus{}wetenschap\PYZus{}schuldenaar\PYZus{}tegenbewijs}
  \PY{k+kr}{consequence} \PY{k+kr}{fulfilled}

  \PY{k+kr}{rule} \PY{n+nv}{vermoeden\PYZus{}wetenschap\PYZus{}wederpartij}
  \PY{k+kr}{under condition}
    \PY{n+nv}{vermoeden\PYZus{}van\PYZus{}toepassing}
    \PY{o}{and} \PY{o}{not} \PY{n+nv}{vermoeden\PYZus{}wetenschap\PYZus{}wederpartij\PYZus{}tegenbewijs}
  \PY{k+kr}{consequence} \PY{k+kr}{fulfilled}

\PY{k+kr}{scope} \PY{n+nc}{FaillissementspaulianaArt42}\PY{o}{:}
\PY{c+c1}{  \PYZsh{} Invoer van de subscope}
  \PY{k+kr}{definition} \PY{n+nv}{bewijsvermoeden\PYZus{}art43}\PY{o}{.}handeling \PY{k+kr}{equals} \PY{n+nv}{handeling}
  \PY{k+kr}{definition} \PY{n+nv}{bewijsvermoeden\PYZus{}art43}\PY{o}{.}faillissement \PY{k+kr}{equals} \PY{n+nv}{faillissement}
  \PY{k+kr}{definition} \PY{n+nv}{bewijsvermoeden\PYZus{}art43}\PY{o}{.}tegen\PYZus{}onderwaarde \PY{k+kr}{equals} \PY{n+nv}{tegen\PYZus{}onderwaarde}
  \PY{k+kr}{definition} \PY{n+nv}{bewijsvermoeden\PYZus{}art43}\PY{o}{.}voldoening\PYZus{}of\PYZus{}zekerheid\PYZus{}niet\PYZus{}opeisbare\PYZus{}schuld \PY{k+kr}{equals}
    \PY{n+nv}{voldoening\PYZus{}of\PYZus{}zekerheid\PYZus{}niet\PYZus{}opeisbare\PYZus{}schuld}
  \PY{k+kr}{definition} \PY{n+nv}{bewijsvermoeden\PYZus{}art43}\PY{o}{.}gelieerde\PYZus{}partij \PY{k+kr}{equals} \PY{n+nv}{gelieerde\PYZus{}partij}
  \PY{k+kr}{definition} \PY{n+nv}{bewijsvermoeden\PYZus{}art43}\PY{o}{.}reeds\PYZus{}verplicht\PYZus{}voor\PYZus{}aanvang\PYZus{}termijn \PY{k+kr}{equals}
    \PY{n+nv}{reeds\PYZus{}verplicht\PYZus{}voor\PYZus{}aanvang\PYZus{}termijn}
  \PY{k+kr}{definition} \PY{n+nv}{bewijsvermoeden\PYZus{}art43}\PY{o}{.}vermoeden\PYZus{}wetenschap\PYZus{}schuldenaar\PYZus{}tegenbewijs \PY{k+kr}{equals}
    \PY{n+nv}{vermoeden\PYZus{}wetenschap\PYZus{}schuldenaar\PYZus{}tegenbewijs}
  \PY{k+kr}{definition} \PY{n+nv}{bewijsvermoeden\PYZus{}art43}\PY{o}{.}vermoeden\PYZus{}wetenschap\PYZus{}wederpartij\PYZus{}tegenbewijs \PY{k+kr}{equals}
    \PY{n+nv}{vermoeden\PYZus{}wetenschap\PYZus{}wederpartij\PYZus{}tegenbewijs}

\PY{c+c1}{  \PYZsh{} Art. 43 vervult de wetenschap aan beide zijden.}
  \PY{k+kr}{exception}
  \PY{k+kr}{rule} \PY{n+nv}{wetenschap\PYZus{}schuldenaar}
  \PY{k+kr}{under condition} \PY{n+nv}{bewijsvermoeden\PYZus{}art43}\PY{o}{.}vermoeden\PYZus{}wetenschap\PYZus{}schuldenaar
  \PY{k+kr}{consequence} \PY{k+kr}{fulfilled}

  \PY{k+kr}{exception}
  \PY{k+kr}{rule} \PY{n+nv}{wetenschap\PYZus{}wederpartij}
  \PY{k+kr}{under condition} \PY{n+nv}{bewijsvermoeden\PYZus{}art43}\PY{o}{.}vermoeden\PYZus{}wetenschap\PYZus{}wederpartij
  \PY{k+kr}{consequence} \PY{k+kr}{fulfilled}
```
\end{Verbatim}




\subsection{Artikel 45 Faillissementswet}



In geval van benadeling door een rechtshandeling om niet, die de
schuldenaar heeft verricht binnen één jaar vóór de faillietverklaring,
wordt vermoed dat hij wist of behoorde te weten dat benadeling van de
schuldeisers het gevolg van de rechtshandeling zou zijn.



\begin{tcolorbox}[colframe=OliveGreen, breakable, title=\textcolor{black}{\texttt{Metadata}},title after break=\textcolor{black}{\texttt{Metadata}},before skip=1em, after skip=1em, left=0.6em, right=0.6em]
\begin{Verbatim}[commandchars=\\\{\},numbers=left,firstnumber=309,stepnumber=1,label={\hspace*{\fill}\texttt{faillissementspauliana.catala\_en}},numbersep=9mm]
```catala
\PY{k+kr}{declaration} \PY{k+kr}{scope} \PY{n+nc}{BewijsvermoedenArt45}\PY{o}{:}
  \PY{k+kd}{input}\PY{+w}{ }\PY{n+nv}{handeling} \PY{k+kr}{content} \PY{n+nc}{Handeling}
  \PY{k+kd}{input}\PY{+w}{ }\PY{n+nv}{faillissement} \PY{k+kr}{content} \PY{n+nc}{Faillissement}
  \PY{k+kd}{input}\PY{+w}{ }\PY{n+nv}{vermoeden\PYZus{}wetenschap\PYZus{}schuldenaar\PYZus{}tegenbewijs} \PY{k+kr}{condition}
  \PY{k+kd}{internal}\PY{+w}{ }\PY{n+nv}{binnen\PYZus{}een\PYZus{}jaar\PYZus{}voor\PYZus{}faillietverklaring} \PY{k+kr}{condition}
  \PY{k+kd}{internal}\PY{+w}{ }\PY{n+nv}{vermoeden\PYZus{}van\PYZus{}toepassing} \PY{k+kr}{condition}
  \PY{k+kd}{output}\PY{+w}{ }\PY{n+nv}{vermoeden\PYZus{}wetenschap\PYZus{}schuldenaar} \PY{k+kr}{condition}
```
\end{Verbatim}

\end{tcolorbox}



\begin{Verbatim}[commandchars=\\\{\},numbers=left,firstnumber=319,stepnumber=1,label={\hspace*{\fill}\texttt{faillissementspauliana.catala\_en}}]
```catala
\PY{k+kr}{scope} \PY{n+nc}{BewijsvermoedenArt45}\PY{o}{:}
\PY{c+c1}{  \PYZsh{} Zelfde datumrekening als bij art. 43 Fw.}
  \PY{k+kt}{date} \PY{k+kd}{round} \PY{k+kr}{down}

\PY{c+c1}{  \PYZsh{} \PYZdq{}binnen één jaar vóór de faillietverklaring\PYZdq{}}
  \PY{k+kr}{rule} \PY{n+nv}{binnen\PYZus{}een\PYZus{}jaar\PYZus{}voor\PYZus{}faillietverklaring}
  \PY{k+kr}{under condition}
    \PY{o}{(}\PY{n+nv}{handeling}\PY{o}{.}datum\PYZus{}rechtshandeling
      \PY{o}{\PYZgt{}}\PY{o}{=} \PY{n+nv}{faillissement}\PY{o}{.}datum\PYZus{}faillietverklaring \PY{o}{\PYZhy{}} \PY{l+m+mi}{1} \PY{o}{year}\PY{o}{)}
    \PY{o}{and} \PY{o}{(}\PY{n+nv}{handeling}\PY{o}{.}datum\PYZus{}rechtshandeling
      \PY{o}{\PYZlt{}} \PY{n+nv}{faillissement}\PY{o}{.}datum\PYZus{}faillietverklaring\PY{o}{)}
  \PY{k+kr}{consequence} \PY{k+kr}{fulfilled}

\PY{c+c1}{  \PYZsh{} Grondslag van het vermoeden: een rechtshandeling om niet binnen 1 jaar.}
  \PY{k+kr}{rule} \PY{n+nv}{vermoeden\PYZus{}van\PYZus{}toepassing}
  \PY{k+kr}{under condition}
    \PY{n+nv}{handeling}\PY{o}{.}om\PYZus{}niet
    \PY{o}{and} \PY{n+nv}{binnen\PYZus{}een\PYZus{}jaar\PYZus{}voor\PYZus{}faillietverklaring}
  \PY{k+kr}{consequence} \PY{k+kr}{fulfilled}

\PY{c+c1}{  \PYZsh{} Weerlegbaar vermoeden: is het ontzenuwd (geslaagd tegenbewijs door de}
\PY{c+c1}{  \PYZsh{} wederpartij), dan vervalt het en valt de wetenschap terug op de hoofdregel}
\PY{c+c1}{  \PYZsh{} van art. 42 Fw (curator moet dan wetenschap aantonen zonder vermoeden).}
  \PY{k+kr}{rule} \PY{n+nv}{vermoeden\PYZus{}wetenschap\PYZus{}schuldenaar}
  \PY{k+kr}{under condition}
    \PY{n+nv}{vermoeden\PYZus{}van\PYZus{}toepassing}
    \PY{o}{and} \PY{o}{not} \PY{n+nv}{vermoeden\PYZus{}wetenschap\PYZus{}schuldenaar\PYZus{}tegenbewijs}
  \PY{k+kr}{consequence} \PY{k+kr}{fulfilled}

\PY{k+kr}{scope} \PY{n+nc}{FaillissementspaulianaArt42}\PY{o}{:}
\PY{c+c1}{  \PYZsh{} Invoer van de subscope}
  \PY{k+kr}{definition} \PY{n+nv}{bewijsvermoeden\PYZus{}art45}\PY{o}{.}handeling \PY{k+kr}{equals} \PY{n+nv}{handeling}
  \PY{k+kr}{definition} \PY{n+nv}{bewijsvermoeden\PYZus{}art45}\PY{o}{.}faillissement \PY{k+kr}{equals} \PY{n+nv}{faillissement}
  \PY{k+kr}{definition} \PY{n+nv}{bewijsvermoeden\PYZus{}art45}\PY{o}{.}vermoeden\PYZus{}wetenschap\PYZus{}schuldenaar\PYZus{}tegenbewijs \PY{k+kr}{equals}
    \PY{n+nv}{vermoeden\PYZus{}wetenschap\PYZus{}schuldenaar\PYZus{}tegenbewijs}

\PY{c+c1}{  \PYZsh{} Het bewijsvermoeden van art. 45 Fw kan de wetenschap schuldenaar vervullen}
  \PY{k+kr}{exception}
  \PY{k+kr}{rule} \PY{n+nv}{wetenschap\PYZus{}schuldenaar}
  \PY{k+kr}{under condition} \PY{n+nv}{bewijsvermoeden\PYZus{}art45}\PY{o}{.}vermoeden\PYZus{}wetenschap\PYZus{}schuldenaar
  \PY{k+kr}{consequence} \PY{k+kr}{fulfilled}
```
\end{Verbatim}

